%----------
%	INSULTOS - MACHINE LEARNING
%----------	

In this case, we followed a similar approach as in Sections \ref{sec:ageneral_ml} and \ref{sec:cnegativo_ml}, totalling 588 experiments in the span of seven months (13/12/23 - 09/07/24), and defining 11 different windows, as seen in Figures~\ref{fig:exp_insultos} and ~\ref{fig:exp_insultos_classes}. 

In particular, Figure \ref{fig:exp_insultos} shows the overall trend in weighted average F1 score, and we can see that the global trend line, with an $R^2 = 0.01$, highlighted a constant trend for the changes made in each window. This phenomenon could be due to the fact that in W2 to W6 we obtained relatively good results in terms of weighted average F1-Score, influenced by the class "Genéricos" but it was not consistent across the other two remaining classes, as seen in Table~\ref{tab:window-metrics-insultos} and Figure~\ref{fig:exp_insultos_classes}. However, we appreciated a positive evolution in the final window (W8), achieving the third highest weighted average F1-Score (0.585) but the highest performance across all classes, with F-Scores of 0.699 for "Sexistas/Misóginos", 0.673 for "Genéricos" and 0.344 for "Deseo de Dañar".


\begin{figure}[H]
	\ffigbox[\FBwidth]
	{\caption{Distribution of the ML Experiments for "Insultos" over Time with the Different Windows and Weighted Avg. F1-Score Trend Line}\label{fig:exp_insultos}}
 % [Aquí ponemos como queremos que aparezca en el índice, sin la referencia]{Aquí como queremos que aparezca en el documento, con la referencia}
	{\includegraphics[scale=0.25]{Plantilla_TFG_ingles_2019/imagenes/Experiments in _Insultos_ over time.pdf}}
\end{figure}

\begin{figure}[H]
	\ffigbox[\FBwidth]
	{\caption{Distribution of the ML Experiments over Time for "Insultos", per Class with the Different Windows (Trend Line of W. Avg. F1-Score in Red)}\label{fig:exp_insultos_classes}}
 % [Aquí ponemos como queremos que aparezca en el índice, sin la referencia]{Aquí como queremos que aparezca en el documento, con la referencia}
	{\includegraphics[scale=0.25]{Plantilla_TFG_ingles_2019/imagenes/Experiments in _Insultos_ over time_classes_trendline.pdf}}
\end{figure}


\begin{table}[H]
    \centering
    \begingroup
    \small % Change to desired font size
    \ttabbox[\FBwidth]{
        \caption{INS Average Evaluation Metrics per Window in ML Experiments}
        \label{tab:window-metrics-insultos} % Unique label for the table
    }{
        \begin{tabularx}{\textwidth}{
            X
            >{\centering\arraybackslash}X
            >{\centering\arraybackslash}X
            >{\centering\arraybackslash}X
            >{\centering\arraybackslash}X
            >{\columncolor{gray!15}\centering\arraybackslash}X % Highlight this column
            >{\centering\arraybackslash}X
            >{\centering\arraybackslash}X
            >{\centering\arraybackslash}X
        }
            \hline
            \rowcolor{gray!30} % Adding a gray color to the header row
            \textbf{} & \textbf{Acc. (Std.)} & \textbf{Time (s)} & \textbf{W. Avg. Prec.} & \textbf{W. Avg. Rec.} & \textbf{W. Avg. F1-Score} & \textbf{Sex./Mis. F1-Score} & \textbf{Gen. F1-Score} & \textbf{D. Dañar F1-Score} \\
            \hline
            W1 & 0.396 (0.059) & 63.067 & 0.383 & 0.396 & 0.365 & 0.290 & 0.473 & 0.354 \\
            \hline
            W2 & 0.522 (0.020) & 476.329 & 0.585 & 0.522 & 0.527 & 0.192 & 0.665 & 0.141 \\
            \hline
            W3 & 0.572 (0.047) & 584.217 & 0.586 & 0.572 & 0.550 & 0.214 & 0.702 & 0.093 \\
            \hline
            W4 & 0.646 (0.048) & 737.703 & 0.586 & 0.646 & 0.602 & 0.221 & 0.785 & 0.041 \\
            \hline
            W5 & 0.617 (0.034) & 952.594 & 0.603 & 0.617 & 0.601 & 0.315 & 0.757 & 0.134 \\
            \hline
            W6 & 0.601 (0.027) & 969.658 & 0.560 & 0.601 & 0.572 & 0.262 & 0.749 & 0.067 \\
            \hline
            W7 & 0.478 (0.042) & 305.891 & 0.396 & 0.478 & 0.405 & 0.222 & 0.587 & 0.138 \\
            \hline
            W8 & 0.392 (0.054) & 248.289 & 0.400 & 0.392 & 0.388 & 0.410 & 0.442 & 0.280 \\
            \hline
            W9 & 0.432 (0.096) & 416.933 & 0.421 & 0.432 & 0.417 & 0.418 & 0.523 & 0.238 \\
            \hline
            W10 & 0.430 (0.089) & 432.454 & 0.419 & 0.430 & 0.412 & 0.377 & 0.524 & 0.254 \\
            \hline
            W11 & 0.607 (0.075) & 1157.582 & 0.605 & 0.607 & 0.585 & 0.699 & 0.673 & 0.344 \\
            \hline
            \multicolumn{9}{l}{Source: 'Training.xlsx'} \\
            \hline
        \end{tabularx}
    }
    \endgroup
\end{table}




\begin{itemize}
    \item \textbf{W1: Baseline Establishment:} For the first phase of the experimentation we employed the data without any modifications and introduced traditional embeddings for three different ML models:
        \begin{itemize}
             \item \textbf{Embedding:} Traditional embeddings explained in Section~\ref{text_encoding_ml}.
             \item \textbf{Models:} Random Forest, Logistic Regression and XGBoosting.
        \end{itemize}
   As seen in Table~\ref{tab:window-metrics-insultos}, the mean performance in terms of weighted F1-Score was solid (0.365), with each class achieving F1-Scores above 0.290. This phase established a baseline for evaluating the next modifications.

    \newpage
    
    \item \textbf{W2: Testing of Upsampling Technique:} In this phase, we tested the effectiveness of the upsampling technique, building on the embeddings used in W1 and incorporating additional ML models. From Table~\ref{tab:window-metrics-insultos}, the results demonstrated an improvement over W1, with the lowest standard deviation (0.020) observed during the experimentation and an overall increase in the mean weighted F1-Score to 0.527. This could be due to the "Genéricos" class, which achieved the highest F1-Score of 0.665. However, the performance for the other two classes, "Sexistas/Misóginos" and "Deseo de Dañar", decreased, with F1-Scores of 0.192 and 0.141, respectively.


    \item \textbf{W3: Testing of SMOTE Technique:} We used the same settings as in W1 and W2 but applying the SMOTE technique. The overall accuracy increased to 0.572, and the weighted average F1-Score improved to 0.550. These enhancements were reflected in better performance for all classes except "Deseo de Dañar", where the F1-Score slightly decreased from 0.141 in W2 to 0.093 in W3. 


    \item \textbf{W4: Testing of ADASYN Technique:} As well as in previous windows, we tested the ADASYN technique with the same configurations. While the accuracy slightly increased to 0.646 and the weighted average F1-Score improved to 0.602, the performance across individual classes was varied. The "Genéricos" class achieved a notable F1-Score of 0.785, the highest observed so far. However, the "Deseo de Dañar" class experienced a decline in performance, with its F1-Score dropping to just 0.041, the lowest recorded across all windows.

    
    \item \textbf{W5: Introduction of Advanced Embeddings and Preprocessing:} In this window we introduced advanced embeddings (Section~\ref{text_encoding_dl}) and some preprocessing to the dataset while testing with the SMOTE technique:
        \begin{itemize}
            \item \textbf{Embedding:} Advanced embeddings were added as introduced in Section~\ref{text_encoding_dl}.
            \item \textbf{Data:} Implemented some text preprocessing, as included in Section~\ref{sec:data_preprocessing}.
        \end{itemize}
    As shown in Table~\ref{tab:window-metrics-insultos}, the accuracy (0.617) and the weighted average F1-Score (0.601) saw a slight decrease compared to W4. Despite this, the "Genéricos" class continued to perform well, achieving an F1-Score of 0.757. Moreover, the "Deseo de Dañar" class showed some improvement with an F1-Score increase to 0.134, though it still was behind other classes.


    \item \textbf{W6: Tokenization:} This phase included tokenization (Section~\ref{special_tokens}) and the use of the ADASYN technique. In general, the mean performance metrics for this window decreased, with an accuracy of 0.601 and an overall weighted F1-Score dropping to 0.572. While the "Genéricos" class performed well with an F1-Score of 0.749, the "Sexistas/Misóginos" and "Deseo de Dañar" classes struggled more, achieving F1-Scores of 0.262 and 0.067, respectively.

    \newpage
    
    \item \textbf{W7: Testing All Modifications with No balancing technique:} In this phase the settings were the same as in W6 but testing without applying any data balancing technique. As expected, the absence of balancing led to a decline in overall performance. The mean weighted F1-Score dropped to 0.405, with the individual class F1-Scores also reflecting this decrease: 0.222 for "Sexistas/Misóginos", 0.587 for "Genéricos", and 0.138 for "Deseo de Dañar". 
    

    \item \textbf{W8: Testing All Modifications with Upsampling technique:} This window was dedicated to testing all previously introduced modifications with the upsampling technique. The results indicated a slight decline in overall accuracy (0.392) and weighted average F1-Score (0.388). However, improvements were observed in the "Sexistas/Misóginos" and "Deseo de Dañar" classes, with F1-Scores rising to 0.410 (compared to 0.222) and 0.280 (compared to 0.138), respectively. 


    \item \textbf{W9: Testing All Modifications with SMOTE technique:} We tested all modifications with the SMOTE balancing technique. The overall performance improved, with a weighted F1-Score of 0.417 but revealed a slight decline (0.238) for "Deseo de Dañar". Despite this, there were improvements in the "Sexistas/Misóginos" and "Genéricos" categories, which saw F1-Scores increase to 0.418 and 0.523, respectively.


    \item \textbf{W10: Testing All Modifications with ADASYN technique:} Same setting as in W7 and W8 but applying ADASYN technique. During this phase, we achieved better performance than in W7 with a weighted F1-Score of 0.412, and across "Genéricos" (0.524) and "Deseo de Dañar" (0.254), except for "Sexistas/Misóginos" (0.377).


    \item \textbf{W11: Introduction of OpenAI Embedding:} In the final phase, OpenAI embedding "text-embedding-3-large" was introduced (Section~\ref{text_encoding_dl}). This phase achieved the best and most balanced overall mean performance across all classes, with an accuracy of 0.607 and a weighted F1-Score of 0.585. The 'Sexistas/misóginos' class had a major improvement and it had the highest F1-Score of 0.699, 'Genéricos' obtained 0.673 while 'Deseo de Dañar' improved to 0.344.
    
    
\end{itemize}



As a result of this analysis, in Table~\ref{tab:insultos-bestmodels} we have collected the top 5 best overall models with their correspondent parameters, and in Table~\ref{tab:insultos-bestmodels_metrics} we show the performance metrics for these models.

\begin{table}[H]
    \centering
    \begingroup
    \small % Change to desired font size
    \ttabbox[\FBwidth]{
        \caption{INS Best ML Models Nomenclature}
        \label{tab:insultos-bestmodels} % Unique label for the table
    }{
        \begin{tabularx}{\textwidth}{
            X  % Name
            X  % Model
            X  % Embedding
            >{\hsize=0.5\hsize\centering\arraybackslas}X  % Embed. Size
            >{\hsize=0.3\hsize\centering\arraybackslas}X  % CV
            X  % Balance
            }
            \hline
            \rowcolor{gray!30} % Adding a gray color to the header row
            \textbf{} & \textbf{Model} & \textbf{Embedding} & \textbf{Embed. Size} & \textbf{CV} & \textbf{Balance} \\
            \hline
            \multicolumn{6}{l}{\textbf{Naive Models}} \\
            \hline
            \insmajor & - & - & - & - & - \\
            \hline
            \insrandom & - & - & - & - & - \\
            \hline
            \multicolumn{6}{l}{\textbf{Finetuned Models}} \\
            \hline
            \insmlptea & MLP & text-embedding-3-large & - & 5 & ADASYN \\
            \hline
            \inslrtesm & Logistic Regression & text-embedding-3-large & - & 5 & SMOTE \\
            \hline
            \insxgbrosm & XGBoosting & Robertuito & 500 & 5 & SMOTE \\
            \hline
            \insrfrosm & Random Forest & Robertuito & 500 & 5 & SMOTE \\
            \hline
            \insxgbbesm & XGBoosting & Beto & 500 & 5 & SMOTE \\
            \hline
            \multicolumn{6}{l}{Source: 'Training.xlsx'} \\
            \hline
        \end{tabularx}
    }
    \endgroup
\end{table}



\begin{table}[H]
    \centering
    \begingroup
    \small % Change to desired font size
    \ttabbox[\FBwidth]{
        \caption{INS Evaluation Metrics for the Best ML Models}
        \label{tab:insultos-bestmodels_metrics} % Unique label for the table
    }{
        \begin{tabularx}{\textwidth}{
            X
            >{\centering\arraybackslash}X
            >{\columncolor{gray!15}\centering\arraybackslash}X % Highlight this column
            >{\centering\arraybackslash}X
            >{\centering\arraybackslash}X
            >{\centering\arraybackslash}X
        }
            \hline
            \rowcolor{gray!30} % Adding a gray color to the header row
            \textbf{} & \textbf{Acc. (Std.)} & \textbf{W. Avg. F1-Score} & \textbf{Sex./Mis. F1-Score} & \textbf{Genéricos F1-Score} & \textbf{D. Dañar F1-Score} \\
            \hline
            \multicolumn{6}{l}{\textbf{Naive Models}} \\
            \hline
            \insmajor &  0.485 (0.035) & 0.318 & 0.000 & 0.653 & 0.000 \\
            \hline
            \insrandom & 0.369 (0.033) & 0.369 & 0.231 & 0.484 & 0.279 \\
            \hline
            \multicolumn{6}{l}{\textbf{Finetuned Models}} \\
            \hline
            \insmlptea & 0.738 (0.047) & 0.740 & 0.889 & 0.744 & 0.609 \\
            \hline
            \inslrtesm & 0.738 (0.061) & 0.732 & 0.842 & 0.773 & 0.571 \\
            \hline
            \insxgbrosm & 0.723 (0.059) & 0.715 & 0.222 & 0.852 & 0.462 \\
            \hline
            \insrfrosm & 0.699 (0.065) & 0.696 & 0.222 & 0.833 & 0.429 \\
            \hline
            \insxgbbesm & 0.723 (0.022) & 0.690 & 0.400 & 0.837 & 0.235 \\
            \hline
            \multicolumn{6}{l}{Source: 'Training.xlsx'} \\
            \hline
        \end{tabularx}
    }
    \endgroup
\end{table}

\newpage

Unlike in Section~\ref{sec:ageneral_ml}, where we were interested in achieving high performance in one specific class ("Comentario Negativo") due to the dependency of subsequent analyses on that particular class, in this case, our objective was to obtain high metrics across all classes. Therefore, we selected model \textbf{\insmlptea} as the best performing model for this task, with a weighted average F1-Score of 0.740 and surpassing 0.600 in all classes. Moreover, it demonstrated the second lowest standard deviation among the models listed in Table~\ref{tab:insultos-bestmodels_metrics}, and significant improvements in all metrics compared to naive models {\insmajor} and {\insrandom}.

Additionally, Figure~\ref{fig:exp_insultos_metrics_best_model} illustrates that the model performed particularly well in the "Genéricos" category with an F1-Score of 0.744 and in the "Sexistas/Misóginos" category with an F1-Score of 0.889. However, the model showed weaker performance in the "Deseo de Dañar" category, achieving a lower F1-Score of 0.609. Specifically, the model frequently misclassified "Deseo de Dañar" and "Sexistas/Misóginos" instances as "Genéricos", which might be attributed to the broad and general nature of the "Genéricos" class.


\begin{figure}[H]
	\ffigbox[\FBwidth]
	{\caption{Confusion Matrix for the Best ML Model (\textbf{\insmlptea}) for "Insultos"}\label{fig:exp_insultos_metrics_best_model}}
 % [Aquí ponemos como queremos que aparezca en el índice, sin la referencia]{Aquí como queremos que aparezca en el documento, con la referencia}
	{\includegraphics[scale=0.15]{Plantilla_TFG_ingles_2019/imagenes/Experiments in _Insultos_ best_model.pdf}}
\end{figure}

\newpage